% 14-cognition.tex: Atlas, Athena, and the layers that must never model a person.
\section{How Polaris observes and understands itself}
\label{sec:cognition}

Two cognitive layers exist, two are proposed, and one line separates all of them from the people
the system serves. The design rule is stated in the repository in one sentence: make power legible,
not people.

\figfile{cognition}{The cognitive layers. Atlas sees what the system is doing; Athena understands
the structure of authority; Metis, proposed, would learn the behaviour of Polaris itself; Myrmex,
proposed in this paper, would attack it from outside. The constitution constrains all of them. The
red line is the person-prohibition: nothing above it models a natural person, and five checks fail
the build if a person-level column or traversal appears in the authority layer.}{fig:cognition}

\subsection{Atlas sees}

The Atlas is the operator's investigation surface: an analytical console that opens on a bounded,
non-geographic overview, a volume time-series with failures overlaid, top-K breakdowns by context,
agency, disclosure and outcome, a two-dimensional pivot, a records grid, and a map kept as a tab for
what a map is good at, one region's drill or one subject's journey. Every view is a bounded
server-side aggregate, never the raw events: at most a few hundred cluster summaries for a viewport,
at most 240 time buckets, at most fifty categories, keyset pagination so that a deep page stays
cheap. A zero-knowledge verification is counted in volume and in the disclosure tallies and is never
located, because it carries no token to attribute it by. The single-entity investigation pages are
built to refuse cross-individual link analysis, and every read of an audit surface is itself
recorded. A live-simulation mode is driven in a headless browser in CI and the console is asserted to
actually stream. \sImpl

\subsection{Athena understands structure}

Athena is a read-only semantic and provenance layer over the authority tables: agencies,
algorithms, contexts, trust attestations and retention policies, as ten object views and eight edge
views, with four functions that answer the governance questions the tree could otherwise answer only
by hand. Why may this agency issue under this algorithm? What proof and disclosure policy bounds
this context? If this algorithm is deprecated, who is affected? Which mechanism enforces this
constitutional rule? The constitution itself is data: C1 to C10 and the vocation are rows, and a
table maps each rule to the exact live mechanism that enforces it, a trigger, a partial unique index,
a named check constraint, a check function or a procedure. A check fails the build if any named
mechanism no longer exists, and a database test proves each is present in the running catalogue,
which closes the drift between prose and code that the constraint lattice alone could not. The
signed registry of Section~\ref{sec:institutions} is a view over Athena, so the map of the
institutions that an outsider fetches is the same map the operator inspects. \sImpl

Athena is safe to exist only because person-legibility is made structurally impossible, not merely
discouraged. Five invariants gate it, each with an adversarial detection test: no Athena view,
function or table references a natural-person table or column, and no column is a stable per-person
surrogate; every function is read-only, so \emph{the graph says revoke, therefore revoked} is
impossible; every function bounds its output; every authority edge resolves to an existing authority
table, and the current views exclude revoked authority, so Athena describes and never owns; and every
rule maps to a mechanism that exists. The same ship that created Athena removed the earlier
person-aggregating views. If the prohibition ever obstructs a feature, the repository's own
assessment says that is the signal to stop, not to relax it. \sImpl

\subsection{What must never be built here}

No person, household or relationship object. No subject handle, opaque or otherwise, stable across
contexts. No traversal from a zero-knowledge event to a subject. No population segmentation or
scoring. No action: the layer reads and never writes. Those are the kill criteria of the design
study that endorsed a constrained Athena, and they are the boundary any future cognitive layer
inherits.

\subsection{Metis and Themis, named carefully}

The roadmap carries a captured brief for a future, advisory-only learning subsystem that would learn
Polaris itself: infrastructure anomaly detection, capacity forecasting with uncertainty, regression
and release-risk intelligence, chaos learning that proposes experiments and never injects them,
root-cause assistance over Athena's dependency graph, and aggregate institutional anomalies of the
form \emph{this agency normally revokes 0.8 percent a month and now revokes 4.9}. Its central
hypothesis is to be falsified before any code: machine learning may infer properties of the system
and must never infer, rank, score or predict the rights of natural persons, and renaming a person
to a subject does not satisfy that if the data stays linkable. Its design law is advisory only:
observe, model, predict, explain, recommend, and never predict then revoke, deny, issue or change a
constitution or a production configuration. The brief proposes the name Metis, because the obvious
alternative collides with the metrics stack, and records that a rule that works as well as a model
should be preferred to the model. Nothing of it is built; the brief calls for a full assessment
first. \sFut

This paper uses the name Themis for the constitution's constraining role over all four layers, and
the repository's own brief uses it in the same sentence: Atlas sees, Athena understands structure,
Metis learns system behaviour, Themis constrains them all. Themis is not a component. It is the
vocation and C1 to C10 as enforced in the schema and pinned by the checks, which already exist. The
name is a reading aid and nothing in the code carries it.

\begin{layercard}{Layer card: system cognition}
\qa{What it is}{Atlas, the bounded operational console; Athena, the read-only authority-and-constitution layer; and two proposals, Metis and Myrmex.}
\qa{Why it exists}{An operator must see what the system is doing and understand where authority comes from; a system this large cannot be governed by hand.}
\qa{What it trusts}{The append-only events and the authority tables it reads.}
\qa{What it refuses}{Any person-level object, handle, traversal, score or action; an unbounded read; a model where a rule would do.}
\qa{What it can see}{Aggregates of events and the structure of authority, keys, contexts, trust and rules.}
\qa{What it cannot see}{A person. The single-entity investigation page exists outside these layers, gated by role and audited on every read.}
\qa{If it fails}{A person reference in Athena fails five checks; an unbounded aggregate fails C8; a learning system that scored people would violate the vocation and is refused before it is designed.}
\qa{Checked by}{\code{check\_athena\_no\_person}, \code{check\_athena\_read\_only}, \code{check\_athena\_functions\_bounded}, \code{check\_athena\_non\_sovereign}, \code{check\_athena\_rule\_enforcement\_resolves}, \code{check\_atlas\_console}, the C8 caps.}
\qa{Now or later}{Atlas and Athena now. Metis after an assessment; Myrmex is a proposal in this paper.}
\qa{Inside and outside}{Inside: operations. Outside: the loop by which the system corrects its own understanding.}
\end{layercard}

\nextlayer{A system that can observe itself still has to learn from what it observes, and the
hardest lesson Polaris has learned is that its own instruments can be green while measuring the
wrong thing. The next section is that lesson.}
